% Library unit: Berger - Kotlerman, Arkhangelsk 1948 (opposite-coloured bishops). Printed once, before
% the first position that uses it; see lib/README.md.

\paragraph{Berger -- Kotlerman, Arkhangelsk 1948}\index{Opposite-coloured bishops!Berger -- Kotlerman}
A classic example of the second defensive system from
Dvoretsky's \emph{Endgame Manual}~\cite{Dvoretsky:2008-endgame-manual} (\xref{Section}{second-defence-system}).
White's bishop stops the b-pawn on the b1--h7 diagonal, and his king keeps watch over the e-pawn.
Every move must be exact.
\diagram[berger-kotlerman]{8/8/8/5B2/1p3b2/2k1p3/8/5K2 w - - 0 1}
1.\K e2 b3 2.\K d1 \K b4 3.\B h7 \K a3 4.\B g6 \K b2 (4...b2, threatening 5...\K a2, 5.\B b1 \K b3
6.\K e2) 5.\B f7! Black threatened 5...\K a1 and then ...b2; White stops this plan by attacking the
b3-pawn. 5...\K a2 6.\B e6 \K a3, threatening 7...b2 8.\B f5 \K a2. 7.\B f5!, and White holds.
